\section{Data}
\textbf{UNIFESP} (Farina, Abdala and Kitamura, Kaggle; CC BY-NC-SA 4.0) holds 2,481 DICOM radiographs. They are de-identified, come from the Universidade Federal de São Paulo under IRB approval, and the Kaggle version is reduced to 256 $\times$ 256 pixels. Radiology residents labelled them on the MD.ai platform, with 20 body regions plus "Others", in a multilabel scheme. The dataset page names four residents; it does not say how many readers saw each image, whether disagreements were adjudicated, or who the patients were. There are 1,738 training images from 768 patients (63 \% of patients have several images) and 743 test images from 330 patients whose labels are hidden. No patient appears in both sets. Ninety-five percent of the images are MONOCHROME1 (inverted). The DICOM \emph{BodyPartExamined} tag is empty in all 2,481 files, so we could not measure the tag error rate on this dataset. Of the training images, 92.4 \% carry one label (122 carry two, 10 carry three), and class frequency runs from 738 (Chest) down to 9 (Clavicles). A random forest on metadata alone (photometric interpretation, bits stored, size; patient-grouped cross-validation) reached accuracy 0.486 against 0.451 for the majority class, with macro-F1 0.059, so acquisition metadata hold almost no label information.
\textbf{Splits.} We split by patient with multilabel stratification (iterative stratification, seed 42). That gave a fixed held-out set of 323 images from 156 patients and a cross-validation pool of 1,415 images from 612 patients in five folds. The held-out set has only 1 to 3 positives for Clavicles (1), Skull (2), Forearm (3), Elbow (3) and Thigh (3), so per-class numbers for these classes are unstable.
\textbf{Other datasets.} For \textbf{MURA} (Rajpurkar et al. 2018) we took the body part from the \texttt{XR\_*} folder name (7 classes) and used the official train/validation split. \textbf{FracAtlas} (Abedeen et al., Sci Data 2023, CC BY 4.0) has 4,083 images from three hospitals in Bangladesh, annotated by two radiologists and an orthopaedist; we used the 3,685 with a single region (hand, hip, leg or shoulder) for zero-shot transfer. \textbf{COVID-19 Radiography} (Chowdhury et al. 2020) has 4 classes and no patient identifiers, so its split is image-level.

\textbf{Dataset availability.} All four datasets are public, but access conditions differ.
\begin{itemize}\item UNIFESP X-Ray Body Part Classification: \url{https://www.kaggle.com/datasets/felipekitamura/unifesp-xray-bodypart-classification} (doi:10.34740/KAGGLE/DSV/3399135; CC BY-NC-SA 4.0).\item MURA (v1.0): Stanford AIMI, \url{https://stanford.redivis.com/datasets/cv1a-apytk3j44} (project page \url{https://stanfordmlgroup.github.io/competitions/mura/}). Access requires AIMI membership and acceptance of the Stanford University School of Medicine Dataset Research Use Agreement (non-commercial, research-only use; redistribution of the data is not permitted). The first author holds approved access under this agreement (status checked on 6 November 2026). Because redistribution is prohibited we do not release MURA images or derived image files, only code, split lists and result files. The experiments reported here used a third-party Kaggle copy of MURA v1.1 (\url{https://www.kaggle.com/datasets/cjinny/mura-v11}, no description, licence unspecified) obtained before the official access was granted \textbf{\textcolor{red}{[VERIFY: compare the md5 hashes of the Kaggle files with the md5\_hash column of the official Redivis table]}}; the official validation split served as our held-out set because the official test set is hidden.\item COVID-19 Radiography Database: \url{https://www.kaggle.com/datasets/tawsifurrahman/covid19-radiography-database} (files (c) original authors; cite Chowdhury et al. 2020 and Rahman et al.).\item FracAtlas: figshare, doi:10.6084/m9.figshare.22363012 (CC BY 4.0). We used the Kaggle copy \url{https://www.kaggle.com/datasets/orvile/fracatlas} with 4,083 images as in the Scientific Data paper; the current figshare version (v7) lists 4,073 images, so the two are not identical.\end{itemize}
No new patient data were collected. Our code, split files, seeds and result files are provided with the submission [repository link on acceptance].

